\documentclass{article}
\title{オプションの数学入門 — ペイオフからグリークスまで}
\author{TeX 図書館}

\begin{document}

\section{この教科書のために — 権利を売買する}

オプションとは「\textbf{将来ある価格で買う(売る)権利}」の取引です。株 100 株を「1 年後に 120 円で買う権利」は、株価が 150 円になれば価値が出て、80 円なら権利は紙屑になります。\textbf{権利の値段}がオプション価格(プレミアム)で、その理論価格を数学で考えるのがこの教科書の主題です。

\begin{quote}
\textbf{【心構え】}\ この教科書は数学の学習用であり、特定の戦略や商品を推奨しません。オプションは損失が急拡大する構造(売り側)を持つため、数理を理解した上で学ぶことが前提です。
\end{quote}

\section{コール・プットとペイオフ}

\textbf{コールオプション}(買う権利)の満期損益は:
\[ \max(S_T - K,\ 0) - C \]
ここで \(S_T\) は満期の株価、\(K\) は行使価格、\(C\) は払ったプレミアム。\textbf{プットオプション}(売る権利)は:
\[ \max(K - S_T,\ 0) - P \]

\begin{quote}
\textbf{【例】}\ 行使価格 100 のコールを 5 円で買った。満期で株価が 120 → 損益 \(+15\)。100 → 損益 \(-5\)(プレミアムだけ失う)。90 → \(-5\)。\textbf{下振れは 5 円で限定、上振れは無限}——買い方の損益の非対称性がオプションの核心です。逆に\textbf{売り側}は利益がプレミアムで限定され、損失が無限になりうる。
\end{quote}

\begin{quote}
\textbf{【演習】}\ 行使価格 100 のプットを 4 円で買ったときの満期損益を、株価 80, 100, 110 のそれぞれで計算せよ。
\end{quote}

\textbf{解答の指針:}\ 80 → \((100-80) - 4 = +16\)。100 → \(-4\)。110 → \(-4\)。プット買いは下振れ保険。

\section{ペイオフ図 — 損益の形}

ペイオフをグラフにすると、ポジションの\textbf{性格}が一目で分かります。

\begin{center}
\begin{tabular}{|l|l|l|l|}
\hline
ポジション & 損益の形 & 最大損失 & 最大利益 \\
\hline
コール買い & 右上がり(肘) & プレミアム & 無限 \\
コール売り & 右下がり(肘) & 無限 & プレミアム \\
プット買い & 左上がり(肘) & プレミアム & 行使価格 − プレミアム \\
プット売り & 左下がり(肘) & 大(株価 0 まで) & プレミアム \\
\hline
\end{tabular}
\end{center}

組み合わせれば形を自由に設計できます。例: \textbf{ストラドル}(コール+プットを同条件で両買い)は「大暴れには当たる、もみ合いは損」という\textbf{変動率への賭け}です。

\begin{quote}
\textbf{【演習】}\ 「株を保有しつつ、行使価格 = 現在株価のプットを買う」(プロテクティブ・プット)のペイオフを述べよ。
\end{quote}

\textbf{解答の指針:}\ 株価の下振れは行使価格で床が張られ(保険金が下りる)、上振れは株として享受できる。コストはプレミアム。株保有 + プット = 同一満期のコール買いの損益と一致する(プット・コール・パリティの入口)。

\section{プレミアムの内訳 — 内在価値と時間価値}

プレミアム \(=\) \textbf{内在価値} \(+\) \textbf{時間価値}。

\begin{itemize}
\item \textbf{内在価値}: 今すぐ権利行使しても得られる額。\(\max(S - K, 0)\)
\item \textbf{時間価値}: 「満期までに株価が動いて価値が出るかもしれない期待」の値。\textbf{OTM(価値なし)でもプレミアムが付く}のはこの分
\end{itemize}

\begin{quote}
\textbf{【例】}\ 株価 105、行使価格 100 のコールのプレミアムが 8 円なら、内在価値 5 + 時間価値 3。\textbf{時間価値は満期で必ず 0 になる}(\( \theta \) の源泉)。
\end{quote}

\begin{quote}
\textbf{【演習】}\ 株価 90、行使価格 100 のコール(OTM)のプレミアムが 3 円とする。内在価値と時間価値に分解せよ。
\end{quote}

\textbf{解答の指針:}\ 内在価値 0、時間価値 3。\textbf{全部が「期待」}であることに注意——買い続けると時間価値の消耗だけで損失が積む(\( \theta < 0 \))。

\section{ボラティリティ — オプション価格の燃料}

オプション価格を左右する最大の変数は\textbf{ボラティリティ}(変動率)です。

\begin{itemize}
\item \textbf{ヒストリカル・ボラティリティ}( HV ): 過去の実際の変動幅。統計で測れる
\item \textbf{インプライド・ボラティリティ}( IV ): 現在の\textbf{オプション価格から逆算した}「市場が想定する将来の変動」
\end{itemize}

IV が上がればコール・プット\textbf{両方}の価格が上がります(変動が大きいほど「権利」の価値が上がるため)。IV は\textbf{イベント前には高く、通過すると低く}なる季節性を持ちます(決算・選挙など)。

\begin{quote}
\textbf{【演習】}\ 「決算前で IV が異常に高い」。この時点でコールを買うことの数理的な注意点を述べよ。
\end{quote}

\textbf{解答の指針:}\ 高 IV はプレミアムに織り込み済み。決算通過で IV が急落すると、株価が正しい方向に動いても\textbf{時間価値の消失が利益を食う}(IV クラッシュ)。買い方は「変動が大きい方向」だけでなく「IV の水準と時期」も問われている。

\section{ブラック・ショールズ式 — 式を読む}

オプション価格の有名な理論式(1973):
\[ C = S_0 N(d_1) - K e^{-rT} N(d_2) \]
\[ d_1 = \frac{\ln(S_0/K) + (r + \sigma^2/2)T}{\sigma \sqrt{T}}, \qquad d_2 = d_1 - \sigma \sqrt{T} \]

式を丸暗記するより\textbf{読み方}を学びます:
\begin{itemize}
\item \(S_0 N(d_1)\): 「株を買う側の期待価値」(\(N\) は標準正規分布の累積確率)
\item \(K e^{-rT} N(d_2)\): 「行使金額を払う側の期待価値」(割り引かれた行使価格)
\item その差が「権利の公平な値段」。**変動 \( \sigma \) が入るのは、株価が触れても届かない世界と、届く世界を区別するため**
\end{itemize}

導出は確率過程(株価を幾何ブラウン運動でモデル化し、リスク中立測度で期待値をとる)に基づきます。『確率過程入門』の二項モデルは、期数を増やした極限でこの式に収束します。

\begin{quote}
\textbf{【演習】}\ BS 式に現れる 5 つの入力(株価・行使価格・満期・金利・ボラティリティ)のうち、\textbf{市場で観測できないもの}はどれか。またその意味を述べよ。
\end{quote}

\textbf{解答の指針:}\ ボラティリティ \( \sigma \) の将来値のみ未知。だから逆に「観測されたオプション価格から \( \sigma \) を求める」(IV)のが実務の使い方。BS は「価格の計算器」でありながら「ボラティリティの測定器」でもある。

\section{グリークス — 感度の 5 文字}

オプション価格が各入力にどう応答するかの指標が\textbf{グリークス}です。

\begin{center}
\begin{tabular}{|l|l|l|}
\hline
グリーク & 定義 & 直感 \\
\hline
\( \Delta \)(デルタ) & \( \partial C / \partial S \) & 株価 1 円動いたら価格はどう動くか(0〜1) \\
\( \Gamma \)(ガンマ) & \( \partial^2 C / \partial S^2 \) & デルタ自身の変化の速さ \\
\( \Theta \)(セータ) & \( \partial C / \partial t \) & 1 日の時間経過で価値が減る速さ \\
\( \mathcal{V} \)(ベガ) & \( \partial C / \partial \sigma \) & IV が 1\%動いたら価格はどうなるか \\
\( \rho \)(ロー) & \( \partial C / \partial r \) & 金利への感度(通常は小さい) \\
\hline
\end{tabular}
\end{center}

\begin{quote}
\textbf{【例】}\ ATM(等価)のコールは \( \Delta \approx 0.5 \)——「半分の確率で行使価値が出る」度合いと読める。ガンマは ATM で最大: 価格が行使価格を跨ぐたびにデルタが 0↔1 を往復するため、\textbf{ATM・満期近}ではセータとガンマが同時に大きくなり(管理が最も難しい地帯)、売り方はこの地帯を避ける慣行があります。
\end{quote}

\begin{quote}
\textbf{【演習】}\ 「コール買い + 株の空売り \( \Delta \) 分」(デルタヘッジ)は、株価の小さな動きに対して中立になる。このポジションが利益を出す唯一の筋合いは何か(ガンマ・セータの観点から)。
\end{quote}

\textbf{解答の指針:}\ 株価の小動きは打ち消されるので、残るのは\textbf{変動(ボラティリティ)}への曝露——実際の変動が IV より大きければ利益(ガンマが働く)。ただし毎日セータで減る。つまり「実現ボラ vs IV」の賭けになっている。

\section{シミュレーションで確かめる}

ブラウン運動で株価を模擬し、オプションのペイオフ分布を見ます。

\begin{lstlisting}[language=Python]
import random, math, statistics

# 幾何ブラウン運動で満期株価を 10000 本生成
S0, sigma, days, trials = 100, 0.20, 90, 10000
dt = 1 / 250
payoffs = []
for _ in range(trials):
    s = S0
    for _ in range(days):
        s *= math.exp(-0.5 * sigma**2 * dt + sigma * random.gauss(0, dt) ** 0.5)
    payoffs.append(max(s - 100, 0))   # 行使価格 100 のコール

print("期待ペイオフ:", round(statistics.mean(payoffs), 2))
print("ペイオフ 0 の割合:", round(
    sum(1 for p in payoffs if p == 0) / len(payoffs), 3))
\end{lstlisting}

\begin{quote}
\textbf{【演習】}\ sigma を 0.20 → 0.40 に変えると「ペイオフ 0 の割合」と「期待ペイオフ」はどうなるか予想してから実行せよ。
\end{quote}

\textbf{解答の指針:}\ 変動が大きくなると 0 で終わる割合は増える一方、当たったときの上振れも大きくなり、\textbf{期待ペイオフは上がる}。コールの期待値は「当たりやすさ」でなく「分布の右裾」で決まる。

\section{おわりに — 権利の値段を考える目}

この教科書の流れ: ペイオフ(§2)で損益の形を掴み、内訳(§3)とボラティリティ(§4)でプレミアムの構造を知り、BS 式(§5)の読み方とグリークス(§6)で感度を測定し、シミュレーション(§7)で分布を自分の手で確かめました。

共通の視点は「\textbf{オプションは株への賭けでなく、変動と時間への賭け}」です。次の一歩は『確率過程入門』(ブラウン運動と二項モデル)と『ベイズ統計入門』(IV の時系列をどう読むか)です。

\begin{quote}
\textbf{【総合演習】}
\begin{enumerate}
\item 内在価値 0 のコールとプットが同条件で同額になるのはどんなときか。
\item デルタ 0.3 のコール 10 枚(1 枚=100 株相当)を買い、デルタヘッジするには株を何株売るか。
\item セータが買い方の敵である理由を時間価値で説明せよ。
\end{enumerate}
\end{quote}

\textbf{解答の指針:}\ 1. 株価 = 行使価格(ATM)のとき。内在価値が両方 0 で、残りが時間価値のみ。2. デルタ合計 \(0.3 \times 1000 = 300\) 株ぶん → 300 株の空売り。3. 時間価値は満期で必ず 0 に向かい、満期が近づくほど消失が加速する(買い方は待つほど不利)。

\end{document}
